% Web専用原稿．
% Lang, Algebra (2002), Ch. XIII §§5–6, pp.522–528 の双線形形式・根基・基底変換を参照．合同変換は係数の変換式として記述し，行列ノートとの相互依存を避ける．
\section{双線形形式}
% genron:begin K-000669
% genron:use K-000001
% genron:use K-000002
% genron:use K-000003
% genron:use K-000004
% genron:use K-000005
% genron:use K-000007
% genron:use K-000017
% genron:use K-000607
% genron:use K-000560
% genron:use K-000608
\begin{defn}[双線形形式]
体 $K$ 上の線形空間 $V,W$ に対し，双線形写像 $b:V\times W\to K$ を双線形形式という．$V=W$ の場合には $V$ 上の双線形形式という．双線形形式全体は点ごとの演算に関して線形空間をなす．
\end{defn}
% genron:end K-000669
% genron:begin K-000670
% genron:use K-000001
% genron:use K-000002
% genron:use K-000003
% genron:use K-000004
% genron:use K-000005
% genron:use K-000007
% genron:use K-000017
% genron:use K-000669
% genron:use K-000562
\begin{defn}[双線形形式の左根基と右根基]
双線形形式 $b:V\times W\to K$ に対し
\[
 L_b=\{v\in V\mid \forall w\in W,\ b(v,w)=0\},\qquad
 R_b=\{w\in W\mid \forall v\in V,\ b(v,w)=0\}
\]
をそれぞれ左根基及び右根基という．双線形性により両者は部分空間である．$L_b=\{0\}$ のとき左非退化，$R_b=\{0\}$ のとき右非退化，双方を満たすとき非退化という．
\end{defn}
% genron:end K-000670
% genron:begin K-000671
% genron:use K-000001
% genron:use K-000002
% genron:use K-000003
% genron:use K-000004
% genron:use K-000005
% genron:use K-000007
% genron:use K-000017
% genron:use K-000670
% genron:use K-000669
% genron:use K-000630
% genron:use K-000631
\begin{prop}[根基による商空間上の非退化形式]
双線形形式 $b:V\times W\to K$ は商線形空間上に非退化双線形形式 $\bar b:(V/L_b)\times(W/R_b)\to K$ を一意に誘導し，
\[
 \bar b(v+L_b,w+R_b)=b(v,w)
\]
を満たす．
\end{prop}
\begin{proof}
$l\in L_b,\ r\in R_b$ に対し $b(v+l,w+r)=b(v,w)$ なので表示は適法である．商空間の演算の定義から双線形性が従う．標準全射の全射性から一意である．$\bar b(v+L_b,w+R_b)=0$ が全ての $w$ について成立すれば $v\in L_b$ であり，左根基は零である．右根基についても同様である．
\end{proof}
% genron:end K-000671
% genron:begin K-000672
% genron:use K-000001
% genron:use K-000002
% genron:use K-000003
% genron:use K-000004
% genron:use K-000005
% genron:use K-000007
% genron:use K-000017
% genron:use K-000669
% genron:use K-000609
% genron:use K-000610
% genron:use K-000600
% genron:use K-000567
\begin{prop}[双線形形式の基底に関する係数表示]
有限基底 $(e_1,\ldots,e_m)$，$(f_1,\ldots,f_n)$ を取る．$g_{ij}=b(e_i,f_j)$ とおけば $x=\sum_i x_ie_i,\ y=\sum_j y_jf_j$ に対して
\[
 b(x,y)=\sum_{i=1}^m\sum_{j=1}^n g_{ij}x_i y_j
\]
である．逆に任意の係数族 $(g_{ij})$ はこの式によって双線形形式を一意に定める．
\end{prop}
\begin{proof}
表示は \GenronRef{K-000609}の二変数の場合である．逆にこの式は基底表示の一意性から適法であり，各変数について和及びスカラー倍を保存する．基底元の組に適用すれば指定係数を得る．一意性は \GenronRef{K-000610}による．
\end{proof}
% genron:end K-000672
\section{対称性と合同変換}
% genron:begin K-000673
% genron:use K-000612
% genron:use K-000001
% genron:use K-000002
% genron:use K-000003
% genron:use K-000004
% genron:use K-000005
% genron:use K-000007
% genron:use K-000017
% genron:use K-000669
% genron:use K-000611
\begin{defn}[対称形式と交代形式]
$V$ 上の双線形形式 $b$ が任意の $x,y\in V$ に対して $b(x,y)=b(y,x)$ を満たすとき対称形式という．任意の $x\in V$ について $b(x,x)=0$ を満たすとき交代形式という．交代形式は $b(y,x)=-b(x,y)$ を満たす（\GenronRef{K-000612}）．
\end{defn}
% genron:end K-000673
% genron:begin K-000674
% genron:use K-000001
% genron:use K-000002
% genron:use K-000003
% genron:use K-000004
% genron:use K-000005
% genron:use K-000007
% genron:use K-000017
% genron:use K-000672
% genron:use K-000673
% genron:use K-000600
% genron:use K-000567
% genron:use K-000609
\begin{prop}[双線形形式の合同変換]
$V$ 上の双線形形式 $b$ と二つの有限基底 $(e_i)_{i=1}^n$，$(e'_j)_{j=1}^n$ を考える．$e'_j=\sum_i p_{ij}e_i$ とし，$g_{il}=b(e_i,e_l)$，$g'_{jk}=b(e'_j,e'_k)$ とおく．このとき
\[
 g'_{jk}=\sum_{i=1}^n\sum_{l=1}^n p_{ij}g_{il}p_{lk}
\]
が成り立つ．係数のこの変換を合同変換という．
\end{prop}
\begin{proof}
$b(e'_j,e'_k)=b(\sum_i p_{ij}e_i,\sum_l p_{lk}e_l)$ を双線形性によって展開する．この公式は基底の変更だけを表すので，形式の対称性・交代性・非退化性は基底の取り方に依存しない．
\end{proof}
% genron:end K-000674
% genron:begin K-000675
% genron:use K-000001
% genron:use K-000002
% genron:use K-000003
% genron:use K-000004
% genron:use K-000005
% genron:use K-000007
% genron:use K-000017
% genron:use K-000673
% genron:use K-000669
% genron:use K-000443
% genron:use K-000612
\begin{prop}[対称形式と交代形式への分解]
$1+1\ne0$ なる体 $K$ 上の線形空間 $V$ の任意の双線形形式 $b$ は，対称形式 $s$ と交代形式 $a$ の和にただ一通り表される．
\end{prop}
\begin{proof}
$2=1+1$ とおき
\[
 s(x,y)=\frac{b(x,y)+b(y,x)}2,\qquad
 a(x,y)=\frac{b(x,y)-b(y,x)}2
\]
と定める．いずれも双線形であり，前者は対称，後者は対角で零である．和は $b$ である．二通りの分解があればその差から対称かつ交代的な形式を得るが，\GenronRef{K-000612}によりその値の二倍は零となる．$2$ は単元なので差は零である．
\end{proof}
% genron:end K-000675
% genron:begin K-000676
% genron:use K-000001
% genron:use K-000002
% genron:use K-000003
% genron:use K-000004
% genron:use K-000005
% genron:use K-000007
% genron:use K-000017
% genron:use K-000673
% genron:use K-000669
% genron:use K-000605
% genron:use K-000604
% genron:use K-000600
% genron:use K-000443
% genron:use K-000100
\begin{prop}[対称双線形形式の直交基底]
$1+1\ne0$ なる体 $K$ 上の有限次元線形空間 $V$ の対称双線形形式 $b$ に対し，相異なる基底元の組について $b(e_i,e_j)=0$ となる基底 $(e_1,\ldots,e_n)$ が存在する．
\end{prop}
\begin{proof}
次元に関する帰納法による．$b=0$ なら任意の基底でよい．$b\ne0$ なら $b(v,v)\ne0$ なる $v$ が存在する．実際，全ての $v$ に対して $b(v,v)=0$ なら $0=b(x+y,x+y)=2b(x,y)$ なので $b=0$ となる．$U=\{x\in V\mid b(v,x)=0\}$ とおくと部分空間である．任意の $x$ に対し
\[
 x=\frac{b(v,x)}{b(v,v)}v+
 \left(x-\frac{b(v,x)}{b(v,v)}v\right)
\]
であり，括弧内は $U$ に属する．$Kv\cap U=\{0\}$ であるから，$U$ の基底に $v$ を加えれば $V$ の基底になる．また $v\notin U$ なので $\dim U<\dim V$ である．帰納法によって $b$ の $U$ への制限について所要の基底を取り，これに $v$ を加える．対称性により $v$ と $U$ の元は双方の順序で直交するから主張が従う．
\end{proof}
% genron:end K-000676
